\bilingualchapter{Framework Details}{框架细节}

\bilingualsection{Rescaling forecasts}{预测值重缩放}

\begin{englishblock} I F YOU'RE GOING TO CREATE YOUR OWN TRADING RULES YOU NEED to rescale them, so that the average absolute value of the forecast is around 10. To do this you need to run a back-test of the trading strategy, although you only require forecast values and you don't need to check performance. You should also average across as many instruments as possible. From the back-test you should measure the average absolute value of forecast values from a back-test, or at least eyeball them to estimate the average absolute forecast. These values should be roughly similar across instruments and long periods of time, otherwise your trading rule could be badly specified and not properly volatility normalised. Once you have the average absolute value then you should divide it into 10. The result is the trading rule's forecast scalar. So for example if the average absolute value was 0.3 then the scalar would be 10 ÷ 0.3 = 33.33. \end{englishblock}

\begin{chineseblock}如果你准备创建自己的交易规则，就需要对这些规则做重缩放，使得预测值的平均绝对值大约在 10 左右。要做到这一点，你需要对该交易策略运行一次回测，不过你只需要预测值本身，并不需要检查策略表现。你还应当尽可能在更多品种之间取平均。通过回测，你应当测量预测值绝对值的平均水平，或者至少通过肉眼估计出平均绝对预测值。这些数值在不同品种和较长时期之间应当大致相似；否则，你的交易规则可能被定义得很糟，没有正确完成波动率归一化。一旦你得到了平均绝对值，就应该用 10 除以它。得到的结果，就是该交易规则的预测缩放因子（forecast scalar）。例如，如果平均绝对值是 0.3，那么缩放因子就是 10 ÷ 0.3 = 33.33。 \end{chineseblock}

\bilingualsection{Calculation of diversification multiplier}{分散化乘数的计算}

\bilingualsubsection{Forecast diversification multiplier}{预测分散化乘数}

\begin{englishblock} Given N trading rule variations with a correlation matrix of forecast values H and forecast weights W summing to 1, the diversification multiplier will be \(\frac{1}{\sqrt{W \times H \times W^T}}\).\footnote{‘T' is the transposition operator. \par “T” 表示转置运算。} Any negative correlations should be floored at zero before the calculation is done, to avoid dangerously inflating the multiplier. \end{englishblock}

\begin{chineseblock}给定 N 个交易规则变体，它们的预测值相关矩阵为 H，预测权重为 W，且这些权重加总为 1，那么分散化乘数为\(\frac{1}{\sqrt{W \times H \times W^T}}\)。在进行计算之前，任何负相关都应当先被截断为 0，以避免把分散化乘数危险地放大。\end{chineseblock}

\bilingualsubsection{Instrument diversification multiplier}{品种分散化乘数}

\begin{englishblock} Given N trading subsystems with a correlation matrix of returns H and instrument weights W summing to 1, the diversification multiplier will be \(\frac{1}{\sqrt{W \times H \times W^T}}\).\footnote{‘T' is the transposition operator. \par “T” 表示转置运算。} Any negative correlations should be floored at zero before the calculation is done, to avoid dangerously inflating the multiplier. \end{englishblock}

\begin{chineseblock}给定 N 个交易子系统，它们的收益相关矩阵为 H，品种权重为 W，且这些权重加总为 1，那么分散化乘数为\(\frac{1}{\sqrt{W \times H \times W^T}}\)。在计算之前，任何负相关都应当先被截断为 0，以避免危险地夸大乘数。\end{chineseblock}

\bilingualsubsection{Spreadsheet example (for either application)}{电子表格示例（适用于以上两种情况）}

\begin{englishblock} For a three asset portfolio, if the correlation matrix is in cells A1:C3 and the relevant weights are in cells F1:F3, then the diversification multiplier will be: \end{englishblock}

\begin{chineseblock}对于一个三资产组合，如果相关矩阵位于单元格 A1:C3，而对应权重位于单元格 F1:F3，那么分散化乘数为：\end{chineseblock}

\begin{quote}
\begingroup\color{\englishbodycolorname} \ttfamily\detokenize{1/SQRT(MMULT(TRANSPOSE(F1:F3), MMULT(A1:C3,F1:F3)))} \endgroup
\end{quote}

\bilingualsection{Calculating price volatility from historic data}{根据历史数据计算价格波动率}

\begin{englishblock} In a spreadsheet package, assuming that the column A contains daily prices, then you first populate column B with percentage returns: \end{englishblock}

\begin{chineseblock}在电子表格软件中，假设 A 列包含每日价格，那么你首先要在 B 列中填入百分比收益：\end{chineseblock}

\begin{quote}
\begingroup\color{\englishbodycolorname} \ttfamily\detokenize{B2 = (A2 - A1) / A1, B3 = (A3 - A2) / A2, ...} \endgroup
\end{quote}

\begin{englishblock} You can then calculate the price volatility from row 26 onwards, using the default of a 25 day moving average: \end{englishblock}

\begin{chineseblock}然后你可以从第 26 行开始，使用默认的 25 日移动平均窗口来计算价格波动率：\end{chineseblock}

\begin{quote}
\begingroup\color{\englishbodycolorname} \ttfamily\detokenize{C26 = STDEV(B2:B26), C27 = STDEV(B3:B27), ...} \endgroup
\end{quote}

\begin{englishblock} The alternative is to use an exponentially weighted moving average (EWMA) of volatility. In general for some variable X if you have yesterday's EWMA Et-1 then today's \end{englishblock}

\begin{chineseblock}另一种做法，是使用波动率的指数加权移动平均（EWMA）。一般来说，对于某个变量 X，如果你已经有了昨天的 EWMA 值 Et-1，那么今天的\end{chineseblock}

\bipairgap

\begin{englishblock} EWMA given a smoothing parameter A is: \end{englishblock}

\begin{chineseblock} EWMA 在给定平滑参数 A 时为：\end{chineseblock}

\[E_t = (A \times X_t) + \left[E_{t-1}(1-A)\right] \]

\begin{englishblock} First of all you need to calculate the A parameter, based on your volatility look-back, using the formula A = 2 ÷ (1 + L). For my suggested default look-back of 36 days, equivalent to a simple moving average of 25 days, we get A = 0.054.\footnote{This is set to give the same half-life as the default look-back of 25 days for a standard moving average. \par 这样设置，是为了让它与默认 25 日简单移动平均具有相同的半衰期。} Assuming you've put 0.054 into cell AA1, and the returns are in column B, in column C we get the squared returns: \end{englishblock}

\begin{chineseblock}首先，你需要根据自己的波动率回看窗口，使用公式 A = 2 ÷ (1 + L) 来计算参数 A。按照我建议的默认回看期 36 天（它等价于 25 天简单移动平均），我们得到 A = 0.054。假设你已经把 0.054 放进单元格 AA1，而收益数据位于 B 列，那么在 C 列中，我们可以计算收益的平方：\end{chineseblock}

\begin{quote}
\begingroup\color{\englishbodycolorname} \ttfamily\detokenize{C2 = B2 ^ 2, C3 = B3 ^ 2, ...} \endgroup
\end{quote}

\begin{englishblock} You set your first estimate of the variance equal to the first square return: \end{englishblock}

\begin{chineseblock}你把第一个方差估计值设为第一个平方收益：\end{chineseblock}

\begin{quote}
\begingroup\color{\englishbodycolorname} \ttfamily\detokenize{D2 = C2} \endgroup
\end{quote}

\begin{englishblock} After that you set the estimate recursively based on your smoothing parameter: \end{englishblock}

\begin{chineseblock}之后，你就根据平滑参数递归地更新这个估计值：\end{chineseblock}

\begin{quote}
\begingroup\color{\englishbodycolorname} \ttfamily\detokenize{D3 = C3 * AA1 + ((1 - AA1) * D2)} \ttfamily\detokenize{D4 = C4 * AA1 + ((1 - AA1) * D3) ...} \endgroup
\end{quote}

\bilingualsubsection{Finally the actual volatility is the square root:}{最后，实际波动率就是它的平方根：}

\begin{quote}
\begingroup\color{\englishbodycolorname} \ttfamily\detokenize{E2 = SQRT(D2), E3 = SQRT(D3), ...} \endgroup
\end{quote}
